\section{Two-point Wick--Ore bars and regular collision data}
\label{wo:sec}

A meromorphic change of chain coordinates can move a contraction from
the differential into the coproduct and the regular coefficient modules.
The homogeneous family below makes this distinction explicit.  With
singleton states fixed, its two-point homology cuts detect the ordered
exchange.  The symmetric Wick contraction cancels from the cycle
subspace.  The contraction
remains in a specified homology marking and in the connecting map of
the regular-extension triangle.

Throughout this section, \(k\) is a field of characteristic zero.
State expressions are finite polynomials, every state generator has
cohomological degree zero, and the internal differential is zero.
The tensor product used here is defined on algebraic configuration
modules.  A comparison with factorization right \(D\)-modules requires
an additional functor preserving the stated operations.

\subsection{The coefficient category and its separated tensor}

An object of \(\mathsf{Fin}_{\mathrm{sep}}\) is a finite set \(I\)
with an equivalence relation.  Its equivalence classes are called
collision classes.  A map \(f:I\to J\) is admissible if
\[
 f(i)\sim_J f(i')\quad\Longrightarrow\quad i\sim_I i'.
\]
Thus labels in distinct source classes remain in distinct target
classes.  Labels in one class may merge or separate.  Disjoint union
places the classes of its two inputs apart; its unit is the empty set.
Let \(cI\) denote the same set with the indiscrete relation.  The
identity on labels gives an admissible refinement \(cI\to I\).

The coefficient algebra is the subalgebra
\begin{equation}\label{wo:eq:coefficients}
 S(I)=k[(z_i-z_j)^{\pm1}:i\not\sim_I j]
       \subset k(z_i:i\in I).
\end{equation}
The additive relations among differences are retained.  In particular,
\(S(cI)=k\).  An admissible map acts by
\(z_i-z_j\mapsto z_{f(i)}-z_{f(j)}\); permitted denominators stay
invertible.  Write \(\partial_i=\partial/\partial z_i\).
For \(s\in S(I)\), the substitution satisfies
\[
 \partial_j f_*(s)
 =f_*\left(\sum_{i:f(i)=j}\partial_i s\right).
\]
This follows on a difference and its inverse, then on the whole
algebra by the Leibniz rule.

A differential configuration module \(M\) assigns a complex of
\(S(I)\)-modules to each \(I\), commuting degree-zero connections
\(\nabla_i\), and semilinear maps \(f_*:M(I)\to M(J)\).
Besides functoriality and compatibility with the differential, the
relations are
\begin{align*}
 \nabla_i(sm)&=(\partial_i s)m+s\nabla_i m,\\
 \nabla_j f_*(m)&=f_*\left(\sum_{i:f(i)=j}\nabla_i m\right).
\end{align*}
Morphisms preserve these data.  Equivalently, choose a small skeleton
and form the \(k\)-linear category \(\mathscr D\) generated by the
coefficients, connections and admissible maps with these relations.
Differential configuration modules are complexes of left
\(\mathscr D\)-modules.  Kernels and cokernels are pointwise.

On \(\mathscr D\), disjoint union uses the inclusions
\(S(I)\otimes S(J)\to S(I\sqcup J)\), the disjoint families of
connections, and disjoint unions of admissible maps.  The defining
relations are preserved by these assignments.  It gives the separated tensor
\begin{equation}\label{wo:eq:tensor}
 (M\odot N)(K)=
 \int^{I,J}\mathscr D(I\sqcup J,K)
       \otimes_k M(I)\otimes_k N(J).
\end{equation}
The coend imposes the relations moving a source operator or source
map between the three factors.  Both bracketings of a triple tensor
are the coend with kernel
\(\mathscr D(I\sqcup J\sqcup H,K)\): composition gives the
identification, and the identity morphism gives its inverse modulo
the coend relations.  The four-variable version proves the pentagon.
Permuting disjoint inputs, with the Koszul sign on complexes, gives
the symmetry.  The representable at the empty set is the unit \(S\).
Indeed, every target connection kills the bare empty map, and then
acts on its coefficient multiples by ordinary differentiation.
The unit is therefore the coefficient module \(S\), with its actual
connections.

\begin{lemma}[Class support]\label{wo:lem:support}
If \(M_1,\ldots,M_n\) vanish at the empty configuration and \(K\)
has \(r\) collision classes, then
\[
 (M_1\odot\cdots\odot M_n)(K)=0\qquad(n>r).
\]
Every admissible map from a disjoint union to \(K\) factors through
the disjoint full unions of target classes met by its inputs.
\end{lemma}
\begin{proof}
Different source classes have images in different target classes.
Each nonzero reduced input has at least one source class, so a
generator of the coend requires at least \(n\) target classes.
Coefficient and connection operators do not change this requirement.
For the factorization assertion, let \(S_a\) be the classes met by
the \(a\)-th input.  These sets are pairwise disjoint.  The map
from that input factors uniquely through the full union \(K_{S_a}\).
The induced map is admissible, proving the assertion.
\end{proof}

Here support refers to the occupied collision classes.  It does not
refer to the support of a sheaf on a collision diagonal.

\subsection{The homogeneous algebra and all-arity associativity}

Let
\begin{equation}\label{wo:eq:states}
 P=k[\tau,\eta,u,v,J_r^+,J_r^-\mid r\ge0],
 \qquad W=\ker(P\xrightarrow{\varepsilon}k),
\end{equation}
where \(\varepsilon\) kills every generator.  The auxiliary
commutative product is the product of normal symbols.  Define its
translation derivation by
\[
 T\tau=T\eta=Tu=Tv=0,\qquad
 TJ_r^\pm=(r+1)J_{r+1}^\pm.
\]
The states \(\tau\) and \(\eta\) remain polynomial generators.
They have not been specialized to nonzero constants.

Set
\[
 \mathcal A(I)=S(I)\otimes_k P^{\otimes I},\qquad
 \nabla_i=\partial_i+T_i.
\]
An admissible map multiplies the normal symbols over each fibre and
inserts the unit at an empty fibre.  Associativity and commutativity
of the auxiliary product give functoriality.  Its derivation rule
and the coefficient chain rule give the connection identities.
Thus \(\mathcal A\) is a differential configuration module.
Its augmentation to \(S\) is the tensor of the maps
\(\varepsilon\); write \(\mathcal I=\ker(\mathcal A\to S)\).
In particular \(\mathcal I(\varnothing)=0\).

On the state factor at label \(i\), define
\[
 U_i=\eta_i\partial_{u_i},\qquad
 V_i=\eta_i\partial_{v_i},\qquad
 D_i=\sum_{r\ge0}
 \left(J_{i,r}^+\partial_{J_{i,r}^+}
       -J_{i,r}^-\partial_{J_{i,r}^-}\right).
\]
These are commuting derivations and commute with \(T_i\).
The sum is finite on every polynomial.  For a set of input labels
\(I\), use \(U_I=\sum_{i\in I}U_i\), and similarly \(V_I,D_I\).
For two disjoint input sets put
\begin{equation}\label{wo:eq:ore}
 O_{I,J}=U_I D_J-V_I U_J,\qquad
 \Theta_{I,J}=e^{O_{I,J}}.
\end{equation}
Each term lowers total \(u,v\)-degree, so the exponential is finite
on a polynomial.  All ordered operators on labelled inputs commute,
because they are polynomials in the commuting \(U_i,V_i,D_i\).

Write \(\partial_{i,r}^\pm=\partial/\partial J_{i,r}^\pm\) and
\begin{equation}\label{wo:eq:jet-kernel}
 c_{rs}(z)=(-1)^r\frac{(r+s+1)!}{r!s!}z^{-r-s-2}.
\end{equation}
The contraction between the two inputs is
\begin{equation}\label{wo:eq:wick}
 \begin{split}
 C_{I,J}={}&\sum_{i\in I,\ j\in J}\tau_i\tau_j
       \sum_{r,s\ge0}c_{rs}(z_i-z_j)\\
 &\hspace{2em}\cdot
 \left(\partial_{i,r}^+\partial_{j,s}^-
       +\partial_{i,r}^-\partial_{j,s}^+\right),
 \qquad E_{I,J}=e^{C_{I,J}}.
 \end{split}
\end{equation}
Only finitely many derivatives act on a polynomial.  Each contraction
lowers the number of currents by two, so this exponential is also
finite.  Both inverse exponentials are obtained by changing the
sign of the exponent.

Let \(m_{I,J}\) concatenate the state tensors and extend their
coefficients to \(S(I\sqcup J)\).  Define
\begin{equation}\label{wo:eq:product}
 \mu_{I,J}=m_{I,J}E_{I,J}\Theta_{I,J}.
\end{equation}
The ordered operator acts first.  There is no commutation assumption
between \(C_{I,J}\) and \(O_{I,J}\).

\begin{proposition}[The separated Wick--Ore algebra]
\label{wo:prop:algebra}
The operations \eqref{wo:eq:product}, the unit \(S\to\mathcal A\)
and the stated augmentation make \(\mathcal A\) an augmented
associative algebra for \(\odot\).  On \(n\) ordered input sets,
every parenthesization is auxiliary concatenation after
\begin{equation}\label{wo:eq:all-arity}
 F_n=\exp\left(\sum_{a<b}C_{I_a,I_b}\right)
       \exp\left(\sum_{a<b}O_{I_a,I_b}\right).
\end{equation}
Each \(F_n\) and \(F_n^{-1}\) preserves the augmentation support
of each participating state factor.
\end{proposition}
\begin{proof}
Every denominator in \(C_{I,J}\) separates the two inputs, so the
output has permitted coefficients.  Naturality under a map within
an input follows from the Leibniz rule: differentiating the product
over a fibre is the sum of the corresponding derivatives before
the merge.  The factors \(\eta_i\) and \(\tau_i\) become the
corresponding central factors at the image label.  An inserted unit
is annihilated by every state derivation.

For connection compatibility, the coefficients satisfy
\[
 \partial_z c_{rs}=(r+1)c_{r+1,s},\qquad
 -\partial_z c_{rs}=(s+1)c_{r,s+1}.
\]
Also \([T_i,\partial_{i,r}^\pm]=-r\partial_{i,r-1}^\pm\),
with zero right side at \(r=0\).  Reindexing the contraction gives
\([T_i,C_{I,J}]=-\partial_i C_{I,J}\).  Hence
\(\nabla_i\) commutes with \(C_{I,J}\), and it commutes with
\(O_{I,J}\) as well.  This proves covariance for every jet.

All contractions commute: their coefficients are independent of
the current variables differentiated by another contraction.
Moreover, \(U_i,V_i\) commute with every contraction, and
\[
 [D_I+D_J,C_{I,J}]=0.
\]
The last identity follows because each summand removes one positive
and one negative current.  An individual charge operator on one
side need not commute with the contraction.

Induct on a parenthesized ordered product.  At its last multiplication,
the inputs split into two consecutive groups \(L,R\).  The outer
ordered operator is \(U_LD_R-V_LU_R\).  It commutes with internal
contractions in \(L\), where it acts only through \(U_L,V_L\),
and with internal contractions in \(R\), where \(D_R\) is their
total charge.  All contractions can therefore be placed to the left
of all ordered operators.  Within each family the operators commute,
and their sums contain every pair exactly once.  This is
\eqref{wo:eq:all-arity}, independent of the tree.

An empty input contributes no derivations, giving the unit.
Give every state generator normal-symbol weight one.  An ordered
correction replaces \(u\) or \(v\) by \(\eta\); a contraction
replaces one current in each factor by its \(\tau\).  These
operations preserve the weight of each state factor.  They therefore
preserve its augmentation support, as do the inverse operators.
This also proves compatibility with the augmentation.
\end{proof}

\subsection{The bar, its signs, and its meromorphic coordinates}

Let \(s\mathcal I\) place a state in degree \(-1\).  The reduced
bar and its coaugmentation are
\[
 \overline B=\bigoplus_{n\ge1}(s\mathcal I)^{\odot n},\qquad
 B=S\oplus\overline B.
\]
Every word is finite.  Lemma~\ref{wo:lem:support} gives a stronger
bound at a fixed configuration: its class arity bounds the word
length, independently of the number of current modes.
We use the positive first multiplication face
\begin{equation}\label{wo:eq:bar-sign}
 b[a_1|\cdots|a_n]
 =\sum_{j=1}^{n-1}(-1)^{j-1}
 [a_1|\cdots|\mu(a_j,a_{j+1})|\cdots|a_n].
\end{equation}
The endpoint faces vanish because every input lies in the
augmentation ideal.  The empty word is \(S\), with zero differential.
Associativity cancels the two mergers of three consecutive inputs;
mergers at disjoint positions cancel with opposite signs.  Hence
\(b^2=0\).  Deconcatenation cuts the ordered word, including its two
unit cuts, and the same accounting proves that it is a chain coproduct.

On these zero-internal-differential states, the convention
\eqref{sb:eq:bar-signs} has the opposite multiplication differential.
Multiplication of a length-\(n\) word by \((-1)^n\) is a chain
isomorphism between the two conventions.  It fixes the empty word
and respects deconcatenation because \((-1)^{p+q}=(-1)^p(-1)^q\).
All formulas in this section, including the connecting map below,
use \eqref{wo:eq:bar-sign}.

Let \(B_0\) be the bar of the same configuration module with its
auxiliary product.  The operators \(F_n\) of
\eqref{wo:eq:all-arity}, together with \(F_0=\id_S\), define a
chain isomorphism
\begin{equation}\label{wo:eq:bar-intertwiner}
 F:B\longrightarrow B_0.
\end{equation}
To prove this, let \(d_i\) multiply adjacent inputs.  Their internal
pair is supplied by \(d_i\), and every other pair by \(F_{n-1}\).
The same total-charge commutation used in
Proposition~\ref{wo:prop:algebra} gives
\(F_{n-1}d_i=d_i^0F_n\).  Summing with
\eqref{wo:eq:bar-sign} proves the chain equation.  Naturality and
connection compatibility descend to the coend, and the inverse
operators are finite on each input.

The standard coproduct need not commute with \(F\).  Its transport
to \(B_0\), on a cut after \(p\) of \(p+q\) factors, is
\begin{equation}\label{wo:eq:transported-cut}
 \Delta^F_{p,q}=(F_p\odot F_q)\Delta_{p,q}F_{p+q}^{-1}.
\end{equation}
Conjugating the chain and coassociativity equations proves those
equations for \(\Delta^F\).  Thus the construction gives a bar
isomorphism with transported operations in every arity.  It does
not identify the transported coproduct with the auxiliary standard
coproduct or assert preservation of regular collision coefficients.

\subsection{The exact two-point complex and its exchange graph}

Take the two separated singleton classes
\(Q=\{1\}:\{2\}\), and write
\[
 z=z_1-z_2,\qquad S_Q=k[z,z^{-1}],\qquad
 M=S_Q\otimes_k W\otimes_k W.
\]
Let \(W_1^{S_Q}=S_Q\otimes(W\otimes1)\) and
\(W_2^{S_Q}=S_Q\otimes(1\otimes W)\).  Then
\(\mathcal I(Q)=W_1^{S_Q}\oplus W_2^{S_Q}\oplus M\).
The only nonzero reduced bar terms are
\begin{equation}\label{wo:eq:two-point-complex}
 \overline B^{-2}(Q)=M_{12}\oplus M_{21},\qquad
 \overline B^{-1}(Q)=W_1^{S_Q}\oplus W_2^{S_Q}\oplus M.
\end{equation}
The two copies record the two orders of the bar inputs.  In each
copy we record the unsuspended coefficient in the fixed label order
\((1,2)\).  This does not identify the two summands.
The coaugmented bar additionally has \(B^0(Q)=S_Q\).

For completeness, each of two reduced inputs must meet a target
class, and their target classes are disjoint.  They must therefore
occupy the two singleton classes in one of their two orders.
The full-class factorization in Lemma~\ref{wo:lem:support} identifies
each allocation with the tensor of the two singleton augmentation
ideals, after scalar extension to \(S_Q\).  A generator and its
image at that full-class factorization agree by the coend relations;
conversely, the singleton inclusions give the inverse on this tensor.
No relation changes the order of the two input slots.  A third
reduced input is impossible.  This proves
\eqref{wo:eq:two-point-complex} in the actual tensor
\eqref{wo:eq:tensor}.

Abbreviate
\[
 \begin{gathered}
 O_{12}=U_1D_2-V_1U_2,\qquad O_{21}=U_2D_1-V_2U_1,\\
 \Theta_{12}=e^{O_{12}},\qquad \Theta_{21}=e^{O_{21}},
 \qquad C=C_{\{1\},\{2\}},\quad E=e^C.
 \end{gathered}
\]
The identity \(c_{sr}(-z)=c_{rs}(z)\) makes \(E\) the same
operator in the two input orders.  Each operator preserves \(M\):
a contraction retains \(\tau_i\) at each participating label.
The differential is consequently
\begin{equation}\label{wo:eq:row}
 b(p,q)=E(\Theta_{12}p+\Theta_{21}q)
          \in M\subset\mathcal I(Q).
\end{equation}
Its one-factor output still carries both insertion labels.  It is
not the constant Laurent coefficient of a field at one point.

Define the ordered exchange and its inverse by
\begin{equation}\label{wo:eq:exchange}
 \mathcal R_{12}=e^{O_{21}-O_{12}},\qquad
 G=\Theta_{21}^{-1}\Theta_{12}=\mathcal R_{12}^{-1}.
\end{equation}
The equality uses commutativity of the ordered operators, without
commuting either one past \(C\).  In the fixed label coordinates,
\(\mu_{12}\mathcal R_{12}=\mu_{21}\).

\begin{theorem}[Two-point homology and cuts]\label{wo:thm:homology}
The map
\begin{equation}\label{wo:eq:graph}
 \gamma:M\longrightarrow M_{12}\oplus M_{21},\qquad
 \gamma(a)=(a,-Ga)
\end{equation}
identifies \(M\) with the literal cycle subspace in degree \(-2\).
The reduced homology is
\[
 H^{-2}(\overline B(Q))\cong M,\qquad
 H^{-1}(\overline B(Q))\cong W_1^{S_Q}\oplus W_2^{S_Q},
\]
with all other reduced groups zero.  In the two independent
allocation summands, the induced reduced coproduct is
\begin{equation}\label{wo:eq:homology-cut}
 \overline\Delta_H\gamma(a)=(a,-Ga).
\end{equation}
Replacing \(C\) by zero and retaining the ordered operators leaves
these two-point cycles, homology and cuts unchanged.
\end{theorem}
\begin{proof}
Invertibility of \(E\) turns \(b(p,q)=0\) into
\(\Theta_{12}p+\Theta_{21}q=0\), hence \(q=-Gp\).
There are no degree-\(-3\) chains, so no incoming boundaries.
Moreover,
\[
 b(\Theta_{12}^{-1}E^{-1}x,0)=x\qquad(x\in M).
\]
Thus the image is exactly the mixed summand \(M\), proving the
homology assertion.  On each singleton the reduced bar is just
\(sW\).  The target of the reduced cut at \(Q\) is therefore
\(M_{12}\oplus M_{21}\), in degree \(-2\) with zero differential.
Cutting the two terms of \(\gamma(a)\) gives
\eqref{wo:eq:homology-cut} directly.  No tensor-product cohomology
identification is needed.  None of these formulas contains \(E\),
which proves the final assertion.  The coaugmentation adds
\(S_Q\) in degree zero and its unchanged unit cuts.
\end{proof}

The ordered exchange is still detected when the singleton states
and their allocation summands are fixed.  For the auxiliary product,
the cut image is the graph \(\{(a,-a):a\in M\}\).
An isomorphism intertwining the cut maps and inducing the identity
on these singleton states must identify that graph with
\(\{(a,-Ga):a\in M\}\).  Projection onto the first coordinate
is an isomorphism on both graphs, so their equality forces \(G=1\).

For example, \(a=u_1\otimes J_2^+\) satisfies
\(Ga=a+\eta_1\otimes J_2^+\).  Its cycle is
\begin{equation}\label{wo:eq:charge-cycle}
 [u_1|J_2^+]-[J_2^+|u_1]-[J_2^+|\eta_1].
\end{equation}
The respective differentials are
\(u_1J_2^++\eta_1J_2^+\), \(u_1J_2^+\), and
\(\eta_1J_2^+\), regarded as mixed tensors.  The signed sum is
zero, whereas omitting the final term leaves
\([\eta_1J_2^+]\ne0\).  In particular \(G\ne1\).
Similarly, \(a=v_1\otimes u_2\) satisfies
\(Ga=a-\eta_1\otimes\eta_2\), giving the cycle
\[
 [v_1|u_2]-[u_2|v_1]+[\eta_2|\eta_1].
\]
Its differential is
\((v_1u_2-\eta_1\eta_2)-v_1u_2+\eta_1\eta_2=0\).
Both central states retain their insertion labels.

\subsection{Marked homology and the order of inverse operators}

Write \(c_0(x)=(x,-x)\) for the auxiliary cycle.  On the two
summands in \eqref{wo:eq:two-point-complex}, the bar isomorphism
\eqref{wo:eq:bar-intertwiner} is
\(F_2=(E\Theta_{12},E\Theta_{21})\), while \(F_1=1\).
Consequently its inverse image is
\begin{equation}\label{wo:eq:marked-cycle}
 \begin{split}
 c_F(x)
 &=\bigl(\Theta_{12}^{-1}E^{-1}x,
             -\Theta_{21}^{-1}E^{-1}x\bigr)\\
 &=\gamma(\Theta_{12}^{-1}E^{-1}x).
 \end{split}
\end{equation}
This is the exact change from the specified auxiliary homology
coordinate \(x\) to the graph coordinate in
\eqref{wo:eq:graph}.  The inverse Wick coefficients in
\eqref{wo:eq:transported-cut} are present in this marking; their
presence alone is not an invariant of unmarked two-point homology.

The inverse order is necessary.  On
\(x=u_1J_1^+\otimes J_2^-\), one has
\[
 O_{12}x=-\eta_1J_1^+\otimes J_2^-,\qquad
 O_{12}Cx=0,\qquad
 CO_{12}x=-z^{-2}\eta_1\tau_1\otimes\tau_2.
\]
Every further term relevant to the two exponential compositions
vanishes.  Therefore
\begin{equation}\label{wo:eq:inverse-order}
 (e^{-O_{12}}e^{-C}-e^{-C}e^{-O_{12}})x
 =z^{-2}\eta_1\tau_1\otimes\tau_2\ne0.
\end{equation}

For the simpler pair \(v=J_1^+\otimes J_2^-\), put
\(t=\tau_1\otimes\tau_2\).  Both ordered operators vanish on
\(v,t\), and \(Cv=z^{-2}t\), \(Ct=0\).  Thus
\begin{equation}\label{wo:eq:current-marking}
 \begin{split}
 \gamma(v)&=[J_1^+|J_2^-]-[J_2^-|J_1^+],\\
 c_F(v)&=\gamma(v)-z^{-2}\gamma(t).
 \end{split}
\end{equation}
The first cycle is regular in its coefficients; the two products
have equal poles, which cancel in its differential.  The second
cycle represents the specified auxiliary coordinate \(v\).
The two central terms in \(\overline\Delta_H\gamma(t)\) belong
to different allocation summands and do not cancel.  What removes
\(E\) from \eqref{wo:eq:homology-cut} is the change of coordinates
\eqref{wo:eq:marked-cycle}.

\subsection{Regular cycle lattices and paired jets}

Extend coefficients from \(S_Q\) to
\(K=k((z))\), and set \(R=k[[z]]\).  Put
\[
 M_R=R\otimes_k W\otimes_k W,\qquad
 M_K=K\otimes_k W\otimes_k W.
\]
These are algebraic tensor products: each vector uses finitely many
state monomials, with coefficients in \(R\) or \(K\).
In particular \(M_R\) is a free \(R\)-module of infinite rank;
no series over all state monomials has been added.
The ordered exponentials and \(G\) are automorphisms of \(M_R\),
since they have no spatial poles.  Formula~\eqref{wo:eq:graph} gives
the regular cycle lattice
\begin{equation}\label{wo:eq:cycle-lattice}
 \ker b\cap(M_R\oplus M_R)=\gamma(M_R).
\end{equation}
Indeed, a regular cycle has first coordinate in \(M_R\), and its
second coordinate is then \(-G\) times the first.  The converse
uses regularity of \(G\).

Under the particular homology marking \(F_H\), identified with
its auxiliary coordinate, this lattice becomes
\begin{equation}\label{wo:eq:marked-lattice}
 F_H\bigl(\gamma(M_R)\bigr)=E M_R.
\end{equation}
To see this, compute
\(F_2\gamma(a)=c_0(E\Theta_{12}a)\) and use
\(\Theta_{12}M_R=M_R\).  Thus the regular cycle lattice is
unchanged on setting \(C=0\), while its image under the fixed
marking retains the contraction.

\begin{proposition}[A paired-jet lattice defect]\label{wo:prop:lattice}
For \(r,s\ge0\), set
\[
 \begin{gathered}
 v=J_{1,r}^+\otimes J_{2,s}^-,\qquad
 t=\tau_1\otimes\tau_2,\qquad p=r+s+2,\\
 \alpha=(-1)^r\frac{(r+s+1)!}{r!s!}\in k^\times,
 \qquad L=Rv\oplus Rt,\quad L'=EL.
 \end{gathered}
\]
On \(Kv\oplus Kt\),
\begin{equation}\label{wo:eq:jet-shear}
 Ev=v+\alpha z^{-p}t,\qquad Et=t.
\end{equation}
The two lattices satisfy
\begin{align}
 L\cap L'&=z^pRv\oplus Rt,&
 L+L'&=Rv\oplus z^{-p}Rt,\label{wo:eq:lattice-sum}\\
 L/(L\cap L')&\cong R/z^pR,&
 L'/(L\cap L')&\cong R/z^pR.\label{wo:eq:lattice-quotients}
\end{align}
For the residue of the ordinary evaluation pairing followed by
\(dz\), the annihilator lattice in the dual plane is
\begin{equation}\label{wo:eq:lattice-dual}
 (L')^\perp=Rv^*\oplus R(t^*-\alpha z^{-p}v^*).
\end{equation}
There is no integer \(N\) for which \(z^NEM_R\subseteq M_R\).
\end{proposition}
\begin{proof}
There is exactly one possible current contraction on \(v\), and
none on \(t\), proving \eqref{wo:eq:jet-shear}.  Every element of
\(L'\) has the unique form
\[
 av+(b+\alpha z^{-p}a)t,\qquad a,b\in R.
\]
It lies in \(L\) exactly when \(a\in z^pR\).  This proves
the intersection and both quotient statements.  Subtracting \(v\)
from \(Ev\) supplies \(\alpha z^{-p}t\); since \(\alpha\)
is a unit, this proves the sum formula.

A dual vector \(Av^*+Bt^*\), with \(A,B\in K\), annihilates
every \(R\)-multiple of \(t\) under residue exactly when
\(B\in R\).  Annihilating all \(R\)-multiples of
\(v+\alpha z^{-p}t\) then requires
\(A+\alpha z^{-p}B\in R\).  These two conditions are precisely
\eqref{wo:eq:lattice-dual}.  The finite coefficient defect has the
perfect residue pairing
\[
 R/z^pR\ \otimes_k\ z^{-p}R/R\longrightarrow k.
\]
In the bases \(1,z,\ldots,z^{p-1}\) and
\(z^{-1},\ldots,z^{-p}\), its matrix is the identity.

Finally, choose \(r,s\) with \(p>N\).  The coefficient of
\(t\) in \(z^NEv\) is \(\alpha z^{N-p}\), which is not
regular.  This proves the failure of a uniform bound.  Each
individual polynomial input still has only finitely many contractions
and a finite pole order.
\end{proof}

The plane \(Kv\oplus Kt\) is a coefficient calculation, not a
translation-stable current submodule: translation reaches higher
jets.  Characteristic zero ensures that \(\alpha\) is a nonzero
unit.  The finite shear formula \eqref{wo:eq:jet-shear} itself has
integer coefficients.  If reduced to a field where that integer
vanishes, it gives no defect on the selected plane.  This observation
does not extend the full exponential construction to positive
characteristic.

\subsection{Termwise regularity and the connecting map}

Give \eqref{wo:eq:two-point-complex}, after scalar extension to
\(K\), the specified termwise lattices
\[
 \begin{split}
 \mathcal L^{-2}&=M_R\oplus M_R,\\
 \mathcal L^{-1}&=(R\otimes W\otimes1)
       \oplus(R\otimes1\otimes W)\oplus M_R.
 \end{split}
\]
For the pure ordered differential
\(b_{\mathrm{ord}}(p,q)=\Theta_{12}p+\Theta_{21}q\), they
form a subcomplex.  For the Wick differential they do not.  The
paired jet in Proposition~\ref{wo:prop:lattice} has
\(\Theta_{12}v=v\), hence
\begin{equation}\label{wo:eq:nonregular-face}
 b(v,0)=v+\alpha z^{-p}t\notin\mathcal L^{-1}.
\end{equation}

\begin{corollary}[No termwise regular chain isomorphism]
\label{wo:cor:regular-chain}
There is no \(K\)-linear chain isomorphism from the Wick complex
to the pure ordered complex which identifies these specified
termwise lattices in both directions.
\end{corollary}
\begin{proof}
If its components are \(H^{-2},H^{-1}\), the chain equation gives
\[
 b=(H^{-1})^{-1}b_{\mathrm{ord}}H^{-2}.
\]
The right side would map \(\mathcal L^{-2}\) into
\(\mathcal L^{-1}\), contradicting
\eqref{wo:eq:nonregular-face}.  No condition fixing the singleton
states is needed for this argument.
\end{proof}

The equal homology and equal regular cycle lattices therefore do not
recover the chain object with its termwise regular data.  There is
also a configuration-theoretic way to retain the missing differential.
Let \(r\) restrict configuration modules to indiscrete objects and
define their regular extension by
\[
 (LV)(I)=S(I)\otimes_k V(cI).
\]
Its connection is the sum of coefficient differentiation and the
state connection.  A map \(LV\to M\) is determined by its
restriction to indiscrete configurations: the value of
\(s\otimes v\) must be \(s(\iota_I)_*v\), where
\(\iota_I:cI\to I\) is refinement.  This formula also constructs
the inverse correspondence, so \(L\dashv r\) and \(rL=\id\).
Both functors are exact over \(k\).

Here \(\mathcal I\) is of regular-extension type, since
\[
 \mathcal I(I)=S(I)\otimes_k
      \ker(P^{\otimes I}\longrightarrow k).
\]
On an indiscrete configuration, Lemma~\ref{wo:lem:support} leaves
only the length-one part of the reduced bar.  The counit therefore
identifies \(Lr\overline B\) with the inclusion
\(s\mathcal I\hookrightarrow\overline B\).
The quotient
\(\mathcal D=\overline B/s\mathcal I\) represents its cone:
the projection from the cone to the quotient has contractible kernel,
the cone of the identity on \(s\mathcal I\).  Thus the exact
regular-extension triangle is
\begin{equation}\label{wo:eq:extension-triangle}
 s\mathcal I\longrightarrow\overline B\longrightarrow
 \mathcal D\longrightarrow(s\mathcal I)[1].
\end{equation}
This extension from indiscrete configurations is distinct from the
choice of \(R\)-lattices above.

At \(Q\), the quotient \(\mathcal D(Q)\) is
\(M_{12}\oplus M_{21}\) in degree \(-2\), with zero differential.
The connecting map takes the differential of a lifted quotient cycle.
With \eqref{wo:eq:bar-sign}, it is exactly
\begin{equation}\label{wo:eq:connecting-map}
 \begin{split}
 \delta:H^{-2}(\mathcal D(Q))&\longrightarrow
                      H^{-1}((s\mathcal I)(Q)),\\
 \delta(p,q)&=E(\Theta_{12}p+\Theta_{21}q).
 \end{split}
\end{equation}
The mixed-support part of the long exact sequence is
\begin{equation}\label{wo:eq:exact-sequence}
 0\longrightarrow M\xrightarrow{\gamma}
 M_{12}\oplus M_{21}\xrightarrow{\delta}M\longrightarrow0.
\end{equation}
Its injectivity, kernel and surjectivity were proved in
Theorem~\ref{wo:thm:homology}.  In the full sequence, the cokernel
of \(\delta\) in \(\mathcal I(Q)\) is
\(W_1^{S_Q}\oplus W_2^{S_Q}\).
The coaugmented version of the triangle additionally retains the
identity on \(S\), which cancels from the quotient.

For the paired jet,
\(\delta(v,0)=v+\alpha z^{-p}t\).  Thus the separate cohomology
objects of \eqref{wo:eq:extension-triangle} lose a contraction
which its connecting map retains, with the state coordinates and
regular coefficient modules fixed.  Meromorphic inverses trivialize
the row over \(K\); Corollary~\ref{wo:cor:regular-chain} shows
why that does not give the specified regular chain comparison.

The product used throughout is the associative separated product
\eqref{wo:eq:product}.  The nonassociative constant field coefficient
of Theorem~\ref{sb:thm:heisenberg-associator} remains a different
operation.  The shared normalized jet coefficient in
\eqref{wo:eq:jet-kernel} and \eqref{sb:eq:heisenberg-kernel}
does not identify their coefficient functors or state objects.
Moreover, specializing \(\tau\) or \(\eta\) to a nonzero
constant prevents the displayed scalar augmentation of \(P\)
from descending: for example \(\varepsilon(\tau-a)=-a\ne0\).

A comparison with the geometric collision recipient must specify
the full coefficient functor, its tensor product, units and
augmentation, and its original regular coefficient modules.  It
must intertwine \eqref{wo:eq:connecting-map} and the operations in
higher configuration arities.  The two-point homology calculation
and the meromorphic bar isomorphism supply explicit data for that
problem; they do not by themselves construct a field identification
or an associative chiral duality.
